% ============================================
% Response to Reviewer 2
% ============================================
\section*{Response to Reviewer 2}

We thank the reviewer for their thoughtful and constructive comments, which have significantly improved the conceptual clarity of the manuscript. Below we address each point.

% ============================================
\subsection*{R2-1. Nutrient enrichment versus interaction strength}

\reviewercomment{The central claim of the paper is that increasing nutrient supply strengthens interactions, which in turn drives a shift toward community-level selection. While this is an appealing and intuitive interpretation, it may be overly simplified.

Nutrient enrichment simultaneously affects multiple ecological processes, including:
\begin{itemize}
\item species' carrying capacities
\item metabolic rates
\item environmental modification (e.g.\ pH shifts)
\item the balance of competition and facilitation
\end{itemize}

Indeed, the manuscript itself shows that pH modification by dominant taxa may contribute to winner identity in the nutrient-rich regime, while not serving as a sole explanation for Dominance frequency. This suggests that the observed transition may reflect changes in environmental feedbacks and dominance structure, rather than a simple increase in pairwise interaction strength.

This raises an important conceptual point:

\emph{Is the observed shift in coalescence outcomes driven by increased interaction strength per se, or by other effects of resource enrichment?}

We therefore suggest reframing this result more broadly in terms of changes in \emph{interaction intensity and environmental feedbacks}, rather than interpreting nutrient enrichment as a direct proxy for stronger interactions. This would align well with recent work highlighting the role of environmentally mediated interactions in microbial communities \citep{Goldford2018,Estrela2021}.}

\responselabel
We agree that this is an important conceptual question to clarify. In our work, the gLV interaction-strength parameter is framed as a phenomenological parameter for ecological coupling within the model, while the nutrient manipulation in Fig.~4 changes nutrient availability and is interpreted through its effects on invasion resistance, net interaction effects, and environmental feedbacks. We therefore use failed pairwise invasion frequency as an operational measure of invasion resistance, in order to approximate how net interaction effects change across nutrient conditions, rather than as a direct measurement of one biochemical mechanism or a literal estimate of every pairwise Lotka--Volterra coefficient. Multiple mechanisms, including resource competition, increased metabolic activity, pH-mediated environmental feedbacks, and changes in the competition--facilitation balance, can contribute to the observed nutrient-dependent change in invasion resistance. This framing is consistent with prior work on nutrient-mediated microbial interactions, environmentally mediated community assembly, and nutrient-concentration-dependent changes in microbial community diversity and structure \citep{Ratzke2018,Ratzke2020,Goldford2018,Estrela2021,DalBello2026,DuanPawar2025}.

\reviewercomment{From a consumer--resource perspective, it is also not mathematically self-evident that increasing resource supply monotonically increases effective pairwise interaction strengths. In our recent theoretical work \citep{DuanPawar2025}, we show that when mapping mechanistic consumer--resource dynamics to effective Lotka--Volterra coefficients, the external supply rate can cancel out under broad conditions, provided that resources are not strictly limiting and metabolic strategies remain fixed.

Under such conditions, resource enrichment increases total biomass but does not necessarily increase mean per-capita competition strength. Instead, enrichment restructures the interaction network through changes in equilibrium abundances and environmental feedbacks. This distinction is important, as it suggests that the observed transition may reflect a reorganisation of interaction structure rather than a simple scaling of interaction strength.

A brief discussion of these issues would significantly strengthen the conceptual clarity of the manuscript.}

\responselabel
The apparent discrepancy comes from how the interaction coefficient is defined. The consumer--resource result discussed by Duan, Pawar and colleagues concerns coefficients derived mechanistically from resource dynamics as effective \emph{per-capita} terms in a population-growth equation. In contrast, our gLV coefficients are coarse-grained effective inhibitory coefficients in a per-capita growth equation within a phenomenological abundance model, and our nutrient-gradient experiment uses pairwise invasion assays to quantify empirical invasion resistance. We therefore do not interpret nutrient enrichment as a direct or monotonic increase in those mechanistic \emph{per-capita} Lotka--Volterra coefficients. As the reviewer correctly notes, resource enrichment can increase total biomass without simply scaling the underlying \emph{per-capita} coefficients. Thus, the nutrient-gradient result should be interpreted as an increase in empirical invasion resistance, not as a claim that nutrient supply directly or monotonically scales every effective \emph{per-capita} interaction coefficient.

Results \S2.4 (``Nutrient-dependent invasion resistance and environmental feedbacks shift coalescence outcomes'') now states, in part: \mschange{``Following prior work showing that nutrient concentration can alter microbial competition and that pH-mediated environmental feedbacks can drive bacterial interactions \citep{Hu2022,Ratzke2018,Ratzke2020,Hu2025}, we conducted additional coalescence experiments by removing glucose and urea from the Base medium or increasing their concentrations (Methods). This yielded two additional media conditions (\figref{fig:fig4}A): Nutr$-$ (no added glucose/urea) and Nutr$+$ (high supplementation). To approximate how net interaction effects change across nutrient conditions, we measured invasion resistance using pairwise invasion assays among the 12 most abundant isolates (95:5 initial frequency; Methods, Supplementary Figs.~7--9).''}

% ============================================
\subsection*{R2-2. Alternative explanations for community-level selection}

\reviewercomment{The manuscript interprets dominance of one parental community and correlated persistence of its taxa as evidence for community-level selection. We find this interpretation interesting and potentially insightful, but it is not uniquely diagnostic on its own.

Similar patterns could arise from:
\begin{itemize}
\item shared environmental filtering
\item correlated traits within parental communities
\item environmental modification (e.g.\ pH tolerance)
\end{itemize}

In particular, the nutrient-rich regime appears consistent with strong environmental filtering imposed by a small number of dominant taxa, which may not require invoking selection acting at the level of the community as a unit.

We therefore recommend clarifying the interpretation of ``community-level selection,'' and explicitly discussing alternative mechanisms that could generate similar patterns. This distinction has been discussed in previous work on community coalescence and microbial assembly \citep{Mansour2018,Rillig2015}.}

\responselabel
We thank the reviewer for raising this important point. We agree that Dominance and correlated persistence should not be interpreted as uniquely identifying a single underlying mechanism. We have therefore revised the manuscript to present Dominance as outcome-level evidence for community-level selection, while defining community-level selection operationally rather than as one exclusive causal process. In the Discussion, we now state that community-level selection describes origin-correlated persistence during coalescence and that this operational definition does not by itself identify causal ecological or biochemical mechanisms, which may vary across systems.

We made this distinction explicit in Results \S2.1 (``Community-level selection is prevalent in microbial coalescence''): \mschange{``Crucially, Dominance provides outcome-level evidence for community-level selection: one parental community largely displaces the other while preserving its compositional structure.''}

At the start of the Discussion paragraph on broader ecological settings, we now make this distinction explicit: \mschange{``Accordingly, we use community-level selection as an outcome-level description of origin-correlated persistence during coalescence. This operational definition does not by itself identify the causal ecological or biochemical mechanisms, which may vary across systems.''}

As an additional plausibility check for environmental feedback, we ran a Ratzke--Gore-style pH-feedback model \citep{Ratzke2018}, which produces asymmetric replacement over an intermediate range of feedback strength (Response Fig.~R2-2 and Supplementary Fig.~43). We agree that pH-mediated environmental modification is a plausible mechanism for coalescence outcomes, and we do not claim that the phenomenological gLV model is the unique explanatory framework. One observation from this comparison is that the pH-only model rarely produces Restructuring outcomes ($\sim 1\%$ across feedback strengths; 3/297 simulated events), whereas Restructuring is observed at substantially higher rates both in our experiments ($\sim 8\text{--}31\%$ across nutrient conditions) and in the phenomenological gLV model ($\sim 22\%$ across $\mu$ values). This comparison does not rule out pH feedback; instead, it shows that pH feedback alone is unlikely to explain the full outcome distribution. As discussed in R2-1, pH-mediated environmental effects can be included implicitly in the coarse-grained gLV coefficients as part of the net inhibitory effect of one species on another, although pH itself is not an explicit state variable in the model.

\begin{figure}[H]
\centering
\includegraphics[width=0.88\textwidth]{internal_Q5_phase_pH.pdf}
\caption*{\textbf{Response Fig.~R2-2.} \textbf{Outcome distribution of a Ratzke--Gore-style pH-feedback model across feedback strengths.} Per-event outcomes in the similarity plane (left, three feedback strengths) and stacked outcome fractions (right) for the pH-feedback model. The model produces asymmetric replacement (Dominance) over an intermediate feedback range, supporting pH-mediated environmental modification as a plausible mechanism contributing to coalescence outcomes. Restructuring outcomes are rare across all feedback strengths ($\sim 1\%$; 3/297 simulated events), in contrast to the substantial Restructuring fractions observed experimentally ($\sim 8\text{--}31\%$) and in the phenomenological gLV model ($\sim 22\%$).}
\end{figure}

In addition, several related comments raise alternative theoretical models and complementary controls, which we address in the corresponding response sections: parental-community biomass effects (R1-1), pH-mediated filtering and winner direction (R1-2, plus the R2-2 and Supplementary Fig.~43 pH-feedback model above), the dominant-taxon removal control for the PDI calculation (R1-3), initial pool size and realized richness effects (R1-4 and R3-3), abundance-skew/additive-null geometric effects (R3-2 and R3-3), and facilitation or competition-only model scope (R3-4). These analyses are now organized across Supplementary Notes~1, 3, and 5: metric and null-model controls are in Supplementary Note~1, simulation-scope controls are in Supplementary Note~3, and biological alternative explanations are in Supplementary Note~5. The passive-retention alternatives are most relevant here: the denser parental community does not generally win, and we did not detect a significant overall effect of initial pool size on Dominance frequency (R1-1, R1-4).

\reviewercomment{Relatedly, we suggest clarifying the ecological meaning of ``failed invasion.'' At present it is primarily used as a proxy for interaction strength, but it is also naturally interpretable as invasion resistance (i.e.\ whether a rare invader can establish). Framing this more explicitly in terms of invasion fitness would help connect the experimental results to ecological theory and may provide additional insight into the observed patterns.}

We thank the reviewer for this important clarification. We agree that failed invasion is naturally interpreted as invasion resistance and is connected to invasion fitness. We therefore revised the manuscript to frame the pairwise assay primarily as an empirical readout of invasion resistance, with invasion fitness as the relevant ecological concept. In Fig.~4B, we retained the y-axis as the measured quantity, ``frequency of failed invasions,'' while the panel framing identifies this as an invasion-resistance analysis. Results \S2.4 now states: \mschange{``To approximate how net interaction effects change across nutrient conditions, we measured invasion resistance using pairwise invasion assays among the 12 most abundant isolates (95:5 initial frequency; Methods, Supplementary Figs.~7--9). The fraction of failed invasions increased monotonically with nutrient supply (Nutr$-$: $2 \pm 1\%$; Base: $33 \pm 4\%$; Nutr$+$: $48 \pm 4\%$; mean $\pm$ s.e.m.; \figref{fig:fig4}B), indicating that invasion resistance increases across the nutrient gradient.''}

The Methods now states that \mschange{``The fraction of failed invasions served as an empirical proxy for net interaction intensity and invasion resistance (\figref{fig:fig4}B).''} We distinguish this two-species assay from the pairwise selection correlation metric in R2-4.

% ============================================
\subsection*{R2-3. Continuous measures alongside categorical outcomes}
\reviewercomment{The classification into Dominance, Mixture, and Restructuring is useful and provides a clear framework for summarising outcomes. However, it depends on thresholding a continuous similarity measure.

As noted in the broader literature \citep{Mansour2018}, coalescence outcomes often lie along a continuum rather than discrete categories.

We suggest that the authors:
\begin{itemize}
\item present continuous similarity measures alongside categorical outcomes
\end{itemize}}

\responselabel
We agree that the classification (Dominance, Mixture, Restructuring) depends on boundaries in a continuous similarity space. PDI was already introduced in Figs.~3 and 4 as a continuous measure of how strongly the coalesced community is directionally weighted toward one parental community versus the other. To make the thresholding explicit, we now present PDI together with asymmetry and retention magnitude as continuous auxiliary metrics underlying the discrete outcome classification, as shown in Response Fig.~R2-3a. We added these analyses to the Supplementary Information as separate medium-specific panels: Supplementary Fig.~10 for Base, Supplementary Fig.~11 for Nutr$-$, and Supplementary Fig.~12 for Nutr$+$. These measures show where each coalescence event lies before category boundaries are applied, while the categorical labels provide the compact summary used in the main analyses.

To verify how sensitive our main nutrient-dependent result is to the classification-boundary choices, we varied the PDI and retention boundaries around the manuscript baseline (Response Fig.~R2-3b). Across a broad range of boundary choices, the Dominance fraction remains ordered Nutr$-$ $<$ Base $<$ Nutr$+$, supporting the robustness of the nutrient-dependent trend. We also added this boundary-sensitivity analysis to the Supplementary Information as Supplementary Fig.~37 and refer to it in Supplementary Note~1.

\begin{figure}[H]
\centering
\includegraphics[width=0.9\textwidth]{marginal_distributions_by_medium.pdf}
\caption*{\textbf{Response Fig.~R2-3a.} \textbf{Continuous similarity measures track the same nutrient-dependent shift as the categorical classifications.} Rows show Nutr$-$, Base, and Nutr$+$; columns show Parental Dominance Index (PDI), where values near 0 or 1 indicate strong asymmetry toward one parental community and values near 0.5 indicate balanced parental contributions; parental asymmetry $|2\mathrm{PDI}-1|$; and squared retention magnitude $r^2 = u^2 + v^2$, which measures how much of the coalesced composition is explained by the parental communities.}
\end{figure}

\begin{figure}[H]
\centering
\includegraphics[width=0.95\textwidth]{boundary_sensitivity_by_medium.pdf}
\caption*{\textbf{Response Fig.~R2-3b.} \textbf{Boundary-sensitivity analysis supports the nutrient-dependent Dominance trend.} (A) Dominance fractions after varying the PDI boundary while holding the retention boundary at $r^2 > 0.5$. (B) Dominance fractions after varying the retention boundary while holding the PDI boundary at the manuscript baseline. Dashed lines mark the manuscript baseline. This sensitivity analysis is descriptive.}
\end{figure}

\reviewercomment{In addition, the distinction between mixture and restructuring is not entirely straightforward biologically, particularly because abundance changes contribute to classification.

We suggest that the authors:
\begin{itemize}
\item clarify the biological interpretation of ``restructuring''
\end{itemize}}

We clarify that Restructuring denotes low parental retention: the coalesced community is not well explained by the parental communities or their combination. In this interpretation, Restructuring can arise when cross-community species pairings that were not present during parental assembly generate a new stable composition, especially in complex multi-species communities where abundance shifts and higher-order interactions make the final state poorly described by the parental-community basis.

Results \S2.1 (``Community-level selection is prevalent in microbial coalescence'') now includes the concise Restructuring definition: \mschange{``Restructuring denotes low parental retention: the coalesced community is not well explained by the parental communities or their combination, and may reflect a new stable composition produced by previously unseen cross-community species interactions.''}

\reviewercomment{We also find the comparison between dot product and Jaccard metrics in Extended Data Fig.~2 potentially very informative. The divergence between these metrics in low-nutrient environments may itself provide evidence that species-level processes dominate in that regime, and this could be discussed more explicitly.}

We agree that this comparison is informative. Extended Data Fig.~2 was generated for the Base-medium outcomes, so we discuss the metric divergence in that context. Consistent with our response to Reviewer~1, we now frame the result as a difference in what each metric measures rather than as a broad robustness claim. This pattern is consistent with the fact that the metrics emphasize different aspects of community composition: abundance-weighted parental-similarity metrics are more sensitive to quantitative dominance in relative abundance, whereas Jaccard reflects presence/absence retention and Jensen--Shannon evaluates divergence across the full abundance distribution. Thus, a coalesced community can be skewed toward one parental community in abundance space while still retaining taxa from both parental communities, or while redistributing subdominant taxa in ways that make it less skewed toward one parental community under Jaccard or Jensen--Shannon.

Results \S2.1 now states the mixed metric result directly: \mschange{``Alternative metric analysis revealed that Dominance remained the most frequent outcome under abundance-weighted compositional metrics, including Euclidean distance and Bray--Curtis dissimilarity, but not under Jaccard index or Jensen--Shannon divergence, which emphasize species-identity retention and full-distribution divergence, respectively (Extended Data Fig.~2).''} Extended Data Fig.~2 caption now carries the metric-divergence interpretation: \mschange{``Dominance is recovered by abundance-weighted parental-similarity metrics, including the primary similarity-map classification, Euclidean distance, and Bray--Curtis dissimilarity, but not by Jaccard index or Jensen--Shannon divergence. This divergence reflects the different biological questions asked by each metric: abundance-weighted metrics quantify whether the coalesced community quantitatively resembles one parent more than the other, whereas Jaccard index measures presence/absence retention and Jensen--Shannon divergence measures differences across the full abundance distribution. Thus, Jaccard and Jensen--Shannon are best interpreted as complementary metric choices that emphasize different features of community composition.''}

% ============================================
\subsection*{R2-4. Interpretation of pairwise selection correlation}

\reviewercomment{The pairwise ``selection correlation'' metric is an interesting attempt to quantify lineage-level coherence, but its interpretation remains somewhat unclear.

It appears closely related to co-occurrence patterns, and it is not fully clear whether the observed correlations reflect ecological interactions, shared environmental responses, or methodological effects. We find this analysis potentially insightful, but its conceptual interpretation would benefit from clarification.}

\responselabel
We agree that the pairwise selection correlation needed a clearer interpretation. In practice, it tests, for each coalescence event, whether species pairs from the same parental community share survival outcomes more strongly than species pairs from different parental communities. The comparison is evaluated against a within-event null in which parental-origin labels are shuffled. This controls for generic shared-fate tendencies within an event, so the measured enrichment specifically asks whether shared fate is associated with parental-community origin. We use this enrichment as evidence for community-level selection, while recognizing that the metric does not identify the causal route; direct competition, shared environmental tolerance, and environmental modification can all generate such structure.

Results \S2.2 now defines the metric directly: \mschange{``We quantified this effect using pairwise selection correlation, which tests across coalescence events whether species pairs from the same parental community share survival outcomes (both persist or both go extinct) more strongly than species pairs from different parental communities (Supplementary Information; \figref{fig:fig2}D).''}

\reviewercomment{We suggest adding a short paragraph linking this metric more explicitly to invasion fitness in a gLV framework. In that context, invasion fitness corresponds to the growth rate of a species when rare in a resident community. The pairwise invasion assays can then be interpreted as empirical approximations of invasion fitness in two-species systems, while coalescence outcomes reflect correlations in invasion success across many species.

This would provide a unifying theoretical framework linking pairwise assays, community coalescence, and the notion of community-level selection.}

\responselabel
We agree that this point needs clarification, especially because ``pairwise'' refers to two different analyses in the manuscript. As discussed in R2-2, the 95:5 pairwise invasion assay measures two-species growth outcomes and is used primarily as an empirical approximation of invasion resistance under each nutrient condition. By contrast, the pairwise selection correlation metric is computed from community coalescence outcomes and asks whether species fates are more strongly coupled within parental communities than between parental communities.

In the gLV framework, these are related but distinct readouts of invasion success and persistence. Pairwise invasion assays approximate rare-invader establishment in two-species resident backgrounds, whereas selection correlation asks whether persistence outcomes are correlated across many species after two assembled communities are mixed. We added this conceptual bridge to Supplementary Note~7, together with an auxiliary gLV analysis showing that excess same-parent shared fate increases with the interaction-strength parameter $\mu$ (Pearson $r = 0.870$, $p = 3.2 \times 10^{-8}$; Supplementary Fig.~36).

Finally, we keep the metric-specific interpretation in Results \S2.2 and Supplementary Note~7, while using the Discussion to state the broader scope of the term: \mschange{``Accordingly, we use community-level selection as an outcome-level description of origin-correlated persistence during coalescence. This operational definition does not by itself identify the causal ecological or biochemical mechanisms, which may vary across systems.''}

% ============================================
\subsection*{R2-5. Pre-selection of natural communities}

\reviewercomment{The inclusion of natural communities is an important strength of the study. However, the use of a common stabilisation phase in defined media may introduce pre-selection effects that reduce ecological heterogeneity across communities.

In particular, incubating diverse natural communities in simplified media is known to drive convergence toward a limited set of functional guilds \citep{Goldford2018}. As a result, the ``natural'' communities used here may already be filtered to resemble the synthetic communities, potentially contributing to the similarity of results across systems.

We therefore suggest that the authors:
\begin{itemize}
\item discuss the potential effects of this pre-selection more explicitly
\item clarify the extent to which taxonomic or functional convergence occurs during the stabilisation phase
\end{itemize}

This would help contextualise the generality of the natural-community results.}

\responselabel
We agree this is an important concern. As the reviewer notes, defined-media stabilization can drive taxonomic or functional convergence \citep{Goldford2018}, and we cannot directly quantify pre-to-post taxonomic convergence or functional convergence here because the present experiments collected post-stabilization 16S community profiles rather than pre-enrichment, metagenomic, metabolomic, or trait-resolved measurements. We therefore narrow our claim: the natural sample-derived experiment shows that the nutrient-dependent increase in Dominance also occurs in taxonomically richer, laboratory-stabilized communities derived from natural samples, not that unmanipulated natural ecosystems follow the same rules.

Within the available data, we asked whether stabilized communities still carried source-specific taxonomic signal. Using same-source biological replicates as a comparison baseline, different-source parental communities remained substantially less similar at the ASV level (Jaccard similarity 1.7- to 4.1-fold lower than same-source replicates across the three media; Response Fig.~R2-5b; Supplementary Fig.~22b), with a comparable pattern at the genus level (Response Fig.~R2-5e; Supplementary Fig.~22e). This does not rule out functional convergence, or taxonomic convergence within each source relative to its original inoculum; only a pre-enrichment baseline could test those possibilities directly, and we do not have one. The data indicate that stabilization did not produce complete ASV- or genus-level taxonomic collapse to a single shared endpoint. Coalesced communities also retained parent-specific taxonomic signal, remaining closer to their own parents than to unrelated same-medium parents (all ASV- and genus-level tests $p<10^{-6}$; Response Fig.~R2-5c,f; Supplementary Fig.~22c,f).

We therefore limit the claim to nutrient-dependent Dominance in taxonomically richer laboratory-stabilized communities derived from natural samples, while noting that taxonomic or functional pre-selection during the seven-cycle stabilization may affect generalization to unmanipulated natural ecosystems. This revision also keeps the natural sample-derived community claims appropriately cautious.

\begin{figure}[H]
\centering
\includegraphics[width=0.95\textwidth]{Fig_R2_5_natural_taxonomic_distinctness.pdf}
\caption*{\textbf{Response Fig.~R2-5.} \textbf{Natural sample-derived communities retain ASV- and genus-level distinctness after stabilization and coalescence.} \textbf{a,d}, ASV- and genus-level richness by medium. \textbf{b,e}, Same-source stabilized parental replicates are more similar than different-source parental communities, indicating retained source-specific taxonomic signal after stabilization. \textbf{c,f}, Coalesced communities remain more similar to their own parental communities than to unrelated same-medium parental communities. Analyses use ASVs above the 0.1\% relative-abundance threshold and post-stabilization 16S profiles; they test for complete post-stabilization taxonomic collapse and retained parent-specific signal, but not pre-to-post enrichment convergence or functional convergence.}
\end{figure}

We revised the manuscript to frame this experiment as applying to natural sample-derived communities established through laboratory stabilization, rather than to unfiltered natural ecosystems. We also added the same-source/different-source taxonomic-distinctness analysis to the Supplementary Information as Supplementary Fig.~22, and revised Results \S2.6 to present pre-selection during stabilization as an alternative explanation for similarity between synthetic and natural sample-derived outcomes:
\mschange{``These communities showed higher Restructuring fractions, possibly reflecting greater taxonomic diversity and more complex interaction networks. Because these communities were laboratory-stabilized before coalescence, pre-selection in defined media could have contributed to their similarity to synthetic-community outcomes. However, their higher richness and retained source-specific taxonomic structure indicate that stabilization did not produce complete ASV- or genus-level taxonomic collapse (Supplementary Figs.~19--22 and Supplementary Note~6). Overall, these results suggest that the nutrient-dependent increase in Dominance observed in synthetic consortia is also present in taxonomically richer natural sample-derived communities.''}
We also added Supplementary Note~6 to state explicitly that pre-selection may occur and that the current data cannot quantify taxonomic convergence during stabilization or functional convergence.


% ============================================
\subsection*{R2-6. Frame gLV as phenomenological}

\reviewercomment{The gLV model successfully reproduces the qualitative transition in coalescence outcomes and serves as a useful minimal framework. However, it is limited to purely competitive interactions and does not explicitly capture environmentally mediated effects such as pH modification.

We therefore suggest framing the model more explicitly as a phenomenological rather than mechanistic description of the system.}

\responselabel
We agree that the gLV model should be framed as a phenomenological competitive baseline rather than as a mechanistic description of the biochemical interactions operating in the experiments. Results \S2.2 (``Theoretical model with random interactions reproduces community-level selection'') now states: \mschange{``The gLV model serves as a phenomenological framework to explore how the statistical properties of interaction coefficients shape coalescence outcomes, rather than as a mechanistic model of the specific biochemical interactions (e.g., pH modification, metabolic cross-feeding) underlying our experimental observations.''}

\reviewercomment{In addition, while the robustness analysis across interaction coefficient distributions is appreciated, a brief discussion of the biological interpretation of these distributions would be helpful. In particular, clarifying how the mean interaction strength parameter ($\mu$) relates to underlying ecological processes would improve interpretability.}

\responselabel
In the model, $\mu$ parameterizes the sampling range for competitive interaction coefficients in the gLV model. Larger $\mu$ makes the competitive coefficient ensemble stronger and more heterogeneous, moving the system away from a near-no-interaction regime and creating stronger mismatches between species that assembled together and outsiders. This lets assembly filter species into mutually compatible groups, the ``cohesion without cooperation'' route to Dominance discussed in the revised Discussion and Supplementary Note~3.

We decoupled these scale and heterogeneity effects in Supplementary Note~3 and Supplementary Fig.~42. This sweep shows that both average inhibitory effect and coefficient heterogeneity contribute to the Dominance regime. This is consistent with the baseline random-interaction model, where the uniform distribution $U(0, 2\mu)$ couples coefficient mean and spread.

We also revised the Results, Methods, and Supplementary Methods to avoid presenting $\mu$ as simply a mean interaction strength and to define it instead as the sampling-range parameter for competitive interaction coefficients. Results \S2.2 now states: \mschange{``We fixed $r_i = 1$ and drew off-diagonal competition coefficients from a uniform distribution $\mathbb{U}(0, 2\mu)$, so increasing $\mu$ expands the sampled coefficient range and increases both the mean and spread of competitive effects \citep{Hu2022,Hu2025}.''} Methods now states: \mschange{``We fixed growth rates to 1 and self-interaction coefficients $\alpha_{ii} = 1$, and drew off-diagonal competition coefficients from a uniform distribution $\mathbb{U}(0, 2\mu)$, so $\mu$ sets the sampled coefficient range and therefore controls both the mean and heterogeneity of competitive effects within this phenomenological model.''} Supplementary Methods further defines $\alpha_{ij}$ as an effective per-capita inhibitory term and states that the baseline $\alpha_{ij}\geq 0$ model captures competitive growth inhibition rather than facilitation.


% ============================================
\subsection*{R2 minor comments}

\reviewercomment{Improve the visualisation of pairwise correlation results, as the separation between groups is not visually clear.}

\responselabel
We agree that the previous visualization could be clearer. We revised the Fig.~2D presentation and caption to make the group summaries and null baseline easier to interpret:
\mschange{``Individual dots show per-event pairwise selection correlations; in the simulation panel, 100 stratified events per group are displayed for legibility. Squares with error bars show mean $\pm$ s.e.m.\ across all coalescence events ($n = 1{,}200$ per group in simulations and all valid experimental events in the experimental panel). Gray horizontal lines indicate the random selection baseline from a null model in which parental-affiliation labels are shuffled (see Supplementary Information).''}

\reviewercomment{Specify how pH was measured (e.g.\ continuously or at endpoints), and whether variability across replicates was assessed.}

\responselabel
We clarified the pH-measurement description in Supplementary Methods, ``Optical Density Measurements and Monoculture Characterization,'' which now specifies the measurement timing:
\mschange{``Community-level pH was measured from culture supernatants at the relevant endpoint, rather than continuously during growth.''}
As with composition, OD, and outcome analyses, biological replicates were treated as separate event-level observations throughout.

\reviewercomment{Clarify the biological interpretation of the interaction coefficient distributions used in the theoretical model.}

\responselabel
We clarified this point in Supplementary Methods, ``Lotka--Volterra Simulations.'' The revised text now explains that $\alpha_{ij}$ is an effective coefficient in the per-capita growth-rate equation that can absorb direct and indirect inhibitory effects, but does not identify a specific biochemical mechanism or dynamically model pH.

In the same Supplementary Methods section, we also clarified the meaning of $\mu$ and the distributional assumption: \mschange{``Thus, in the baseline competitive ensemble, $\mu$ sets the range of sampled competitive interaction coefficients. This parameter changes the overall scale and heterogeneity of pairwise coefficients, rather than isolating a single biological mechanism.''}

\reviewercomment{Ensure consistent terminology and notation throughout (e.g.\ ``pairwise'', ``community-level'', $\alpha_{ij}$).}

\responselabel
We standardized the terminology throughout the revised manuscript. In particular, Results \S2.2 now uses ``pairwise selection correlation'' for the species-pair fate-correlation metric and defines it as follows:
\mschange{``We quantified this effect using pairwise selection correlation, which tests across coalescence events whether species pairs from the same parental community share survival outcomes (both persist or both go extinct) more strongly than species pairs from different parental communities (Supplementary Information; \figref{fig:fig2}D).''}
We use ``community-level selection'' for origin-correlated persistence at the parental-community level and use $\alpha_{ij}$ consistently for the gLV interaction coefficients.

\reviewercomment{Correct minor typographical issues, including ``generalis ability'' $\to$ ``generalisability''.}

\responselabel
We corrected the ``generalis ability'' typographical error and addressed other minor typographical issues throughout the revised manuscript.
